\section{Abstract Universal and Universal Cause}

The metaphysical objection to particularity has intuitive force because human minds habitually ascend toward truth by leaving particulars behind. We understand a forest by grasping what oak and pine share; we understand mathematics by stripping away wooden sticks and chalk marks to study pure number and figure; we formulate physical laws by abstracting from the quirks of individual experiments. Because our science grows more powerful as its concepts grow wider, we quietly project our own cognitive ascent onto reality itself. We assume that the higher a being stands, the more general and less particular it must be---and that God, standing at the summit of all things, must be the supreme abstraction.

Classical Catholic metaphysics dismantles that assumption at its root. When St.\ Thomas Aquinas examines how creatures exist, how human minds know, and how God creates, knows, and governs, he exposes a fatal equivocation inside the word ``universal.'' There are two radically different kinds of universality: the \textit{abstract universal} of human predication, and the \textit{universal cause} of real being. The first is a mark of finite cognitive poverty; the second is the plenitude of divine power.

\subsection{The Poverty of Abstraction and the Reality of Singulars}

In extra-mental reality, no abstract universal walks the earth. ``Man in general'' never speaks, suffers, chooses, or exists; only Peter, Mary, Socrates, and this or that concrete human being subsist. In his early treatise \textit{De ente et essentia} (\textit{On Being and Essence}), Aquinas states the point directly: if the common character of a species belonged to human nature simply as it exists outside the mind, then commonality would be found inside Socrates himself; yet that is false, because in Socrates no community is found, and whatever is in him is individuated.\endnote{St.\ Thomas Aquinas, \textit{De ente et essentia}, c.\,3 (\latin{in Socrate non invenitur communitas aliqua, sed quicquid est in eo est individuatum}), ed.\ Ludwig Baur and Josef Koch (M\"unster: Aschendorff, 1933). All English renderings of \textit{De ente et essentia} are close inline translations from this Latin edition.} Universality as a logical property belongs to a nature only as it exists in an intellect that compares one abstracted likeness to many individuals outside the mind.

Why, then, does the human intellect traffic first in universals rather than singulars? Not because universals are richer than concrete realities, nor because singulars are beneath genuine knowledge, but because the human intellect is the lowest rank of intelligence---a spiritual power belonging to a soul that is the substantial form of a material body. In the \textit{Summa Theologiae}, Aquinas explains that our intellect cannot directly grasp a material singular in its concrete individuality; it must receive sensory impressions (phantasms) and strip away individualizing material conditions---``this flesh and these bones''---in order to lay hold of the intelligible form.\endnote{Aquinas, \textit{ST} I, q.\,85, a.\,1, corpus and ad 2 (trans.\ Fathers of the English Dominican Province, Project Gutenberg eBook 17611).} Abstraction is an honest act---it does not lie when it considers a shared nature without asserting that the nature floats apart from matter---because ``the mode of understanding, in one who understands, is not the same as the mode of a thing in existing.''\endnote{Aquinas, \textit{ST} I, q.\,85, a.\,1, ad 1.} Yet abstraction is also a loss of resolution. To know the actual singular---to judge that ``Socrates is a man,'' or to act morally in concrete life---our intellect must bend back indirectly toward the sensory phantasm.\endnote{Aquinas, \textit{ST} I, q.\,86, a.\,1, corpus and ad 2.}

Aquinas also refuses the Platonic inference that singularity itself resists intelligibility: ``Intelligibility is incompatible with the singular not as such, but as material, for nothing can be understood otherwise than immaterially.''\endnote{Aquinas, \textit{ST} I, q.\,86, a.\,1, ad 3.} Abstraction is therefore not an ontological promotion of reality from low particularity to high generality; it is an epistemic compression performed by finite minds that cannot hold all the richness of a concrete creature in one intellectual gaze. To gain breadth across a class, our concepts surrender depth in the individual. The wider an abstract universal grows---from \textit{Socrates} to \textit{man} to \textit{animal} to \textit{body} to \textit{substance} to \textit{being in general}---the thinner its content becomes.

\subsection{God Is Not a Genus and Not Being in General}

If God were universal in the way a logical genus or an abstract concept is universal, God would be the thinnest, most indeterminate notion conceivable---a blank logical container awaiting particular differences to fill it out. Aquinas blocks that confusion completely. First, God is not contained in any genus, neither as a species under a higher class nor as a principle reducible to one category: ``all in one genus agree in the quiddity or essence of the genus... but they differ in their existence... But in God they do not differ.''\endnote{Aquinas, \textit{ST} I, q.\,3, a.\,5, corpus.} Because God's essence is identical with his act of being (\latin{ipsum esse subsistens}, subsistent being itself), he has no potentiality to be narrowed by a specific difference, nor can he share a common category with creatures; rather, ``God is the principle of all being. Therefore He is not contained in any genus as its principle.''\endnote{Aquinas, \textit{ST} I, q.\,3, a.\,5, corpus.}

Second, and even more decisively for the scandal of particularity, Aquinas warns against mistaking God's subsistent being for \latin{esse commune}---the abstract universal ``being in general'' shared conceptually by everything that exists. In \textit{De ente et essentia}, immediately after showing that God is pure \latin{esse}, he cautions that calling God being alone must not lead us into the error of those who claimed that God is that universal being by which each thing formally exists.\endnote{Aquinas, \textit{De ente et essentia}, c.\,5 (\latin{Nec oportet, si dicimus quod Deus est esse tantum, ut in illorum errorem incidamus, qui Deum dixerunt esse illud esse universale, quo quaelibet res formaliter est}).} Abstract common being is indeterminate because it neither includes nor excludes further additions; it can be contracted into stone, oak, angel, or human precisely because it is incomplete. God's being is the exact opposite: nothing can be added to it because it already contains every perfection in pure simplicity, so that by its very purity it is being distinct from all being.\endnote{Aquinas, \textit{De ente et essentia}, c.\,5 (\latin{unde per ipsam suam puritatem est esse distinctum ab omni esse}).}

This distinction overturns the entire metaphysical front of the scandal. A finite universal cannot coincide with concreteness: the more general it is, the less determinate it is. The divine \textit{universal cause} obeys the inverse law. Because God is pure subsistent act, wholly unrestricted by any created genus, his universality is not the emptiness of an abstraction stretched thin over many things, but the inexhaustible fullness of the First Being from whom every finite thing receives its whole act of existence. As Aquinas demonstrates in the treatise on creation, ``all things which are diversified by the diverse participation of being, so as to be more or less perfect, are caused by one First Being, Who possesses being most perfectly.''\endnote{Aquinas, \textit{ST} I, q.\,44, a.\,1, corpus.}

Nor is this universal causality a distant initial push after which creatures drift on their own. Because a creature's essence is distinct from its act of existence, it no more sustains itself in being than the daytime air retains sunlight after the sun sets. ``For the being of every creature depends on God, so that not for a moment could it subsist, but would fall into nothingness were it not kept in being by the operation of the Divine power.''\endnote{Aquinas, \textit{ST} I, q.\,104, a.\,1, corpus.} Creation and conservation are one continuous, timeless bestowal of existence to each concrete thing, here and now.

\subsection{Innermost Presence, Knowledge of Singulars, and Particular Providence}

Three consequences follow immediately from God's universality as cause of being---and each one reverses the philosophical prejudice against the particular.

First, transcendence does not create distance; it enables immediate presence within every creature. Finite agents stand outside one another and push from the surface; the higher a creaturely ruler sits, the farther removed he is from the details of his realm. God is not a finite ruler scaled up. Because created existence (\latin{esse}) is God's proper effect, and because an agent must be immediately joined to that upon which it acts, God is present wherever anything exists: ``God is in all things; not, indeed, as part of their essence, nor as an accident, but as an agent is present to that upon which it works... But being is innermost in each thing and most fundamentally inherent in all things since it is formal in respect of everything found in a thing... Hence it must be that God is in all things, and innermostly.''\endnote{Aquinas, \textit{ST} I, q.\,8, a.\,1, corpus.} Nothing is too small, too material, or too local for the Creator's immediate touch, because nothing exists at all except by his continuous gift of its innermost act of being. As St.\ Paul told the philosophers on the Areopagus, God is ``not far from every one of us. For in him we live and move and are'' (Acts 17:27--28).

Second, God knows every singular directly, down to the last detail of its individuality. Classical Aristotelian commentators and Avicennian thinkers had struggled here: if intellect is immaterial and matter is the principle of individuation, how can the supremely immaterial Mind know this particular sparrow or this particular human grief? Some tried to say God knows singulars only through universal laws, as an astronomer predicts eclipses from planetary orbits without knowing this eclipse in its concrete singularity. Aquinas rejects that compromise as unworthy of the Creator. An astronomer does not create the moon and the sun; a human craftsman produces only the form of a chair out of pre-existing wood. God's creative knowledge causes the whole being of the creature---not merely its shared specific form, but its very matter and every individuating principle:

\begin{quote}
[S]ince God is the cause of things by His knowledge, as stated above (A. 8), His knowledge extends as far as His causality extends. Hence as the active power of God extends not only to forms, which are the source of universality, but also to matter, as we shall prove further on (Q. 44, A. 2), the knowledge of God must extend to singular things, which are individualized by matter. For since He knows things other than Himself by His essence, as being the likeness of things, or as their active principle, His essence must be the sufficing principle of knowing all things made by Him, not only in the universal, but also in the singular.\endnote{Aquinas, \textit{ST} I, q.\,14, a.\,11, corpus.}
\end{quote}

In replying to the first objection, Aquinas states the contrast between human abstraction and divine knowledge plainly: ``the intelligible species in the divine intellect, which is the essence of God, is immaterial not by abstraction, but of itself, being the principle of all the principles which enter into the composition of things, whether principles of the species or principles of the individual; hence by it God knows not only universal, but also singular things.''\endnote{Aquinas, \textit{ST} I, q.\,14, a.\,11, ad 1.}

Third, divine providence reaches every singular creature without crushing secondary causality, contingency, or human freedom. Against ancient views that confined God's care to incorruptible heavens or broad biological species while abandoning perishable individuals to blind chance, Aquinas insists that ``all things are subject to divine providence, not only in general, but even in their own individual selves.''\endnote{Aquinas, \textit{ST} I, q.\,22, a.\,2, corpus.} Far from demeaning a supreme ruler to attend to particulars, ``every operative science is the more perfect, the more it considers the particular things with which its action is concerned.''\endnote{Aquinas, \textit{ST} I, q.\,22, a.\,3, ad 1.}

At the same time, immediate providence over the design of all things does not turn creatures into passive puppets. God executes his governance through created intermediaries---natural causes, angels, human parents, prophets, teachers, and free choices---``not on account of any defect in His power, but by reason of the abundance of His goodness; so that the dignity of causality is imparted even to creatures.''\endnote{Aquinas, \textit{ST} I, q.\,22, a.\,3, corpus.} Because God is the transcendent cause of being itself, he does not compete on the same plane with finite causes: he prepares necessary causes for what happens of necessity and contingent causes for what happens contingently, preserving human freedom under providence without destroying it.\endnote{Aquinas, \textit{ST} I, q.\,22, a.\,4, corpus and ad 1, 3; q.\,22, a.\,2, ad 4.} Nor does particular providence make God the author of sin. He permits particular defects and moral evils without willing or causing the evil of fault, ordering the whole universe so that greater goods---such as the preservation of natural species and the patience of martyrs---may be drawn even from what creatures corrupt.\endnote{Aquinas, \textit{ST} I, q.\,22, a.\,2, ad 2.}

\subsection{Personhood as the Summit of Particularity}

There is one further turn in the metaphysical reversal. Particularity is not merely tolerated in creation; it ascends in rank as being ascends, reaching its highest created dignity in the rational \textit{person}.

In every category of being, individual instances exist, yet substance holds individuality in a primary way because it subsists in itself rather than inhering in another, as color inheres in a wall. Within the order of substances, however, Aquinas identifies a still loftier singularity:

\begin{quote}
Further still, in a more special and perfect way, the particular and the individual are found in the rational substances which have dominion over their own actions; and which are not only made to act, like others; but which can act of themselves; for actions belong to singulars. Therefore also the individuals of the rational nature have a special name even among other substances; and this name is ``person.''\endnote{Aquinas, \textit{ST} I, q.\,29, a.\,1, corpus.}
\end{quote}

A stone or an oak tree is a genuine singular, yet it is moved toward its end by its natural impulses alone. A human being---an ``individual substance of a rational nature,'' in Boethius's classic definition adopted by Aquinas---holds dominion over its own acts through intellect and free will. Because each free act belongs to an irreplaceable ``who'' rather than a fungible instance of a species, the individual is no longer subordinate to the species as a mere disposable specimen. Aquinas therefore concludes: ``\,`Person' signifies what is most perfect in all nature---that is, a subsistent individual of a rational nature.''\endnote{Aquinas, \textit{ST} I, q.\,29, a.\,3, corpus.}

Catholic philosophers and theologians have long debated the precise metaphysical principle by which material individuals of the same species are distinguished from one another. Aquinas locates the principle of individuation in material substances in designated matter (\latin{materia signata}, matter considered under determinate dimensions), while maintaining that the human soul, once individuated by its union with this body, retains its individuated being permanently even after death.\endnote{Aquinas, \textit{De ente et essentia}, c.\,2 and c.\,5.} Bl.\ John Duns Scotus, judging matter too indeterminate to account for the positive perfection of singularity, proposed instead a positive contracting principle---an individualizing reality formally distinct from the common nature, later termed \latin{haecceitas} (``thisness'')---while Francisco Su\'arez and other scholastic masters held that every real entity is individuated by its own concrete entirety. Those accounts remain legitimate school positions rather than defined dogmas of the Church. What unites the Catholic metaphysical tradition across its schools is the substantive verdict: singularity is not a stain upon being or an exile from intelligibility. In the rational person, particularity is the crown of the created cosmos.

This conclusion also preserves strict apophatic reserve. When the Church applies the name ``person'' to God---confessing one divine essence in three distinct Persons, Father, Son, and Holy Spirit---she does so analogically, never univocally. As Aquinas notes in replying to the objection that matter individuates while God is immaterial, ``God cannot be called an `individual' in the sense that His individuality comes from matter; but only in the sense which implies incommunicability.''\endnote{Aquinas, \textit{ST} I, q.\,29, a.\,3, ad 4.} Our speech about God always falls short of his uncreated mystery; we know that he is, and that every pure perfection pre-exists in him in a higher mode, without comprehending the divine essence as it is in itself. Yet that very reserve destroys the rationalist dogma that God must resemble an abstract concept. Because the Creator is neither a finite particular competing against other particulars nor a blank universal hovering above them, there is nothing incoherent---and nothing beneath divine majesty---in his acting through concrete history, choosing particular messengers, and addressing individual persons by name.
